\documentclass[10pt,a4paper]{amsart}
\usepackage[margin=19mm]{geometry}
\usepackage{amsmath,amssymb,mathtools}
\usepackage[scr=rsfs,cal=euler]{mathalfa}
\usepackage{xfrac}
\usepackage[unicode,hidelinks,bookmarks=false]{hyperref}
\newcommand{\ie}{i.e.\ }
\newcommand{\cf}{cf.\ }
\newcommand{\resp}{respectively\ }
\newcommand{\Subsection}{\subsection}
\providecommand{\nobreakdash}{\nobreak\hbox{-}}
\setlength{\emergencystretch}{2em}
\multlinegap0pt
\usepackage{bm}
\usepackage{bbm}
\usepackage{dsfont}
\usepackage{xspace}
\usepackage{cancel}
\usepackage{mathtools}
\usepackage{cleveref}
\usepackage{amsthm}
\makeatletter
\@ifundefined{claim}{\newtheorem{claim}{Claim}}{}
\@ifundefined{lemma}{\newtheorem{lemma}{Lemma}}{}
\@ifundefined{proposition}{\newtheorem{proposition}{Proposition}}{}
\makeatother
\newcommand{\Exp}[1]{\mathbb{E}\left[#1\right]}
\newcommand{\Pro}[1]{\mathbb{P}\left(#1\right)}
\newcommand{\aas}{a.a.s.\xspace}
\newcommand{\bigO}{\mathcal{O}}
\newcommand{\eps}{\varepsilon}
\newcommand{\hrg}{\mathcal{G}(n, \alpha, C)}
\title[Distributed Symmetry Breaking on Hyperbolic Random Graphs]{Distributed Symmetry Breaking on Hyperbolic Random Graphs}
\author{Yannic Maus, Janosch Ruff, Sonia Simons, George Skretas}
\date{Source: arXiv:2607.09170v1 (CC BY 4.0)}
\begin{document}
\maketitle

\noindent\emph{Source and licence.} Distributed Symmetry Breaking on Hyperbolic Random Graphs. Version arXiv:2607.09170v1 (CC BY 4.0), distributed under the Creative Commons Attribution 4.0 International licence, as recorded in the arXiv licence metadata for this exact version and shown on the abstract page https://arxiv.org/abs/2607.09170v1. Full rights notice: licenses/arxiv-2607.09170-CC-BY-4.0.txt.\par\smallskip

\noindent\textit{Editorial scope.} Counted unit: the Proposition whose source label is \texttt{luby-aint-shatter-my-hrg} in the article (source0.tex lines 1916--1917, source label \texttt{luby-aint-shatter-my-hrg}), delivered together with the article's own complete original proof of it (source0.tex lines 1919--1992, one \texttt{proof} environment of 74 lines). No mathematics is written, repaired or re-derived: statement and proof are the article's own text line for line. Required context retained uncounted: lem:luby-obstruction (lines 1828--1838). Every definition, display and label of those lines is retained unchanged. Omitted from this package and not redistributed: 1--1827, 1839--1915, 1918--1918, 1993--2099. Nothing inside a retained excerpt is omitted or altered: every retained body is delivered exactly as the article's lines are written, so no omission receipt of any kind accompanies this package. Citations: the retained excerpts carry no citation command at all, so no bibliography entry of the article is transcribed and no reference is redistributed. Reference adaptation: none was needed and none was made; every reference inside the delivered unit resolves inside it, so no unresolved reference or citation is produced. Printed slips kept literal: no printed word, symbol, hypothesis or number of the counted statement or of its proof was altered. Layout and class: the article's own class is \texttt{article} with packages including geometry, lineno, mathtools, fontenc, fnpct, microtype, paralist, libertine, libertinust1math, sectsty, todonotes, ulem, hyperref, enumerate; the wrapper is the lane's pinned \texttt{amsart} editorial wrapper at 10pt on A4, which restates the theorem-like environments the retained text needs and adds the article's own macro definitions that the retained text needs (\texttt{\textbackslash Exp}, \texttt{\textbackslash Pro}, \texttt{\textbackslash aas}, \texttt{\textbackslash bigO}, \texttt{\textbackslash eps}, \texttt{\textbackslash hrg}), with extra wrapper packages (bm, bbm, dsfont, xspace, cancel, mathtools, cleveref). The delivered numbering is the unit's own. Assets: no retained span contains a figure environment, a graphics-inclusion command, a PDF-page inclusion, a table or a TikZ picture; no image file is copied into this package, the package declares no asset dependency and no image file is redistributed with it. Licence: version arXiv:2607.09170v1 of this article is distributed under the Creative Commons Attribution 4.0 International licence, as recorded in the arXiv licence metadata for this exact version and shown on the abstract page https://arxiv.org/abs/2607.09170v1. A full rights notice is delivered at licenses/arxiv-2607.09170-CC-BY-4.0.txt. Characters: every retained excerpt is pure ASCII, so the delivered engine, XeLaTeX with its default Latin Modern text font, reports no missing character.
\begin{lemma}[Luby Obstruction]\label{lem:luby-obstruction}
    Let $G \sim \hrg$ be a threshold hyperbolic random graph and let $\eps < 1 - 1/(2\alpha)$ be a constant larger than $0$. Then, \aas there exist a set of vertices $U(\eps) \subseteq V$ of size $|U(\eps)| \in \Omega\left(n^{1 - 1/2\alpha -\eps}\right)$, where a vertex $u \in U(\eps)$ has the following properties.
%
    \begin{enumerate}
           \item\label{item:property1}\textbf{Degree of $u$}: The degree of $u$ is $\deg(u) \in \Theta(n^{\eps / 2})$.
        \item\label{item:property2}\textbf{Leaves of $u$}: Vertex $u$ has $\Theta(n^{\eps / 2})$ neighbours that are leaves.
        \item\label{item:property3}\textbf{Similar degree path of $u$}: For any constant $t \in \bigO(1)$, $u$ has a similar degree path $W(u)$ of length at least $2t$.
    \end{enumerate}
     Moreover, for any pair $u, u' \in U(\eps)$ it holds that for any pair $v \in L(u) \cup W(u) \cup \{u\}$ and $v' \in L(u') \cup W(u') \cup \{u'\}$ that $\{v,v'\} \not\in E(G)$.
%
\end{lemma}

\begin{proposition}\label{luby-aint-shatter-my-hrg}  Let $G \sim \hrg$ be a threshold hyperbolic random graph and let $t \in \bigO(1)$ be any fixed constant. Then $\Delta(G_t) \in n^{\Omega(1)}$ \aas   
\end{proposition}

\begin{proof}
With hindsight, we set $\eps := \frac{1 - \frac{1}{2\alpha}}{2t(t+2)}$ and consider the set of vertices $U(\eps)$ of \Cref{lem:luby-obstruction}. We show that after $t$ rounds there exists \aas a vertex $u \in U(\eps)$ such that $\deg_t(u) \in n^{\Omega(1)}$.

To this end, for $i \in [t]$, let $L_{(i)}(u) := \{v \in N(u) \cap V_{(i)} : \deg(v) = 1\}$ (the ``remaining leaves'' of $u$ after round $i$), $W_{(i)}(u):=
\left(\bigcup_{j=1}^{2(t-i)}\{w_j\}\right)\cap V_{(i)}$\footnote{Recall that $W(u) := \{w_1, w_2, \dots w_k\}$ is the similar degree path of $u$ such that for $j \in [k]$ it holds $\{w_{j}, w_{j+1}\} \in E$.} (the ``remaining'' similar degree path of $u$ after round $i$) and we define the random variable
%
\begin{align}
    K_{(i)} := |\{u \in U(\eps) \cap V_{(i)} : |L_{(i)}(u)| \in n^{\Omega(1)} \text{ and } |W_{(i)}(u)| \geq 2(t-i)\}|.  
\end{align}
%
Note that if $K_{(t)} \geq 1$ \aas this is sufficient to prove our desired statement and that $K_{(0)} \in \Omega\left(n^{1 - 1/2\alpha - \frac{1 - \frac{1}{2\alpha}}{2t(t+2)}}\right)$ \aas due to \Cref{lem:luby-obstruction} and our choice of $\eps$. In order to prove $K_{(t)} \geq 1$ \aas, we will use the following claim, which tells us that if $K_{(i)}$ is polynomial in $n$, then so is $K_{(i+1)}$ \aas given that $i < t$.
%
\begin{claim}\label{claim:appendix}
Let $\eps := \frac{1 - \frac{1}{2\alpha}}{2t(t+2)}$. Then, for $i \in [t]$ it holds
%
$$
\Pro{K_{(i+1)} \in \Omega\left(n^{1 - 1/2\alpha - \eps -(i+1)\eps(2t+1)} \right) | K_{(i)} \in  \Omega\left(n^{1 - 1/2\alpha - \eps -i\eps(2t+1)} \right)} \in 1 - o(1).
$$
%
\end{claim}
%
{
\renewcommand{\qedsymbol}{$\blacksquare$} % or any symbol you like
\begin{proof}[Proof of claim] 
Let $U_{(i)} = \{u \in U(\eps) \cap V_{(i)} : |L_{(i)}(u)| \in n^{\Omega(1)} \text{ and } |W_{(i)}(u)| \geq 2(t-1)\}$, i.e., $|U_{(i)}| =  K_{(i)}$. We analyse one round of Luby on $U_{(i)}$. To this end, we fix a vertex $u \in U_{(i)}$ and consider the following events:
%
\begin{enumerate}
    \item[($\mathcal{E}_u$)] Vertex $u$ draws a random number for Luby that lies in the open interval $(1- \frac{2t + 1}{n^\eps}, 1- \frac{2t}{n^\eps})$.
    
    \item[($\mathcal{E}_{\text{path}}$)] For $j \in [2(t-i)+1 ]\setminus\{0\}$, vertex $w_j \in W(u)$\footnote{Where $w_j \in W(u)$ is the $j$-th vertex of the similar degree path $W(u)$ such that $\{u, w_1\} \in E$ and  $\{w_j, w_{j+1}\} \in E$.} draws a random number for Luby that lies in the open interval $(1- \frac{2t + 1 -i}{n^\eps}, 1- \frac{2t- i}{n^\eps})$. 
    
\item[($\mathcal{E}_{\text{leaves}}$)] Any vertex $v \in N((W(u) \cup \{u\}) \setminus (W(u)) \cup \{u\})$  draws a random number for Luby that lies in the open interval $(0, 1- \frac{2t + 1}{n^\eps})$.   
\end{enumerate}
%
Observe that if event $\mathcal{E}_u \cap \mathcal{E}_{\text{path}} \cap \mathcal{E}_{\text{leaves}} =: \mathcal{E}$ occurs, then $u \in U_{(i+1)}$. To see this, note that all leaves of $u$ draw a smaller number than $u$ and $u$ is not removed as it is the largest number among neighbours, except for $w_1 \in N(u)$, which does not join the independent set since $w_2 \in N(w_1)$ draws a larger number than $w_1$. This ``chain'' of events continues for the entire similar degree path $W_{(i)}(u)$ by the intersection of the two events $\mathcal{E}_{\text{path}}$ and $\mathcal{E}_{\text{leaves}}$. Hence, only the "terminal" vertex $w_{2(t-i)}$ of $W_{(i)}(u)$ that might join the independent set, removing all its neighbours including the neighbour $w_{2(t-i-1)}$ but no other vertex of the similar degree path $W_{(i)}(u)$. Since $u$ has polynomially many leaves, all necessary conditions for $u \in U_{(i+1)}$ are fulfilled. We continue by lower bounding the probability of event $\mathcal{E}$.

\paragraph{Event $\mathcal{E}_u$:} Since the range of the open interval is $n^{-\eps}$ we have
%
\begin{align}\label{eq:randomdraw1}
    \Pro{\mathcal{E}_u} = n^{-\eps}.
\end{align}
%
\paragraph{Event $\mathcal{E}_{\text{path}}$:} Similar the event $\mathcal{E}_u$, the range of the interval of each random number is $n^{-\eps}$. Since each draw is independent and we draw at most $2t$ random numbers, we obtain
%
\begin{align}\label{eq:randomdraw2}
    \Pro{\mathcal{E}_{\text{path}}} \geq n^{-2t\eps}.
\end{align}
%
\paragraph{Event $\mathcal{E}_{\text{leaves}}$:} Since for any $w \in W(u) \cup \{u\}$ we have $\deg(w) \in \bigO(n^{\eps/2})$, we can upper bound $m := |N((W(u) \cup \{u\}) \setminus (W(u)) \cup \{u\})| \in \bigO(n^{\eps/2})$ using that $|W(u)|$ is constant since $t \in \bigO(1)$. Applying this with the fact that the desired event for every vertex $v$ of drawing a random number within the interval $(0, 1- \frac{2t + 1}{n^\eps})$ has probability $1- \frac{2t + 1}{n^\eps}$, we obtain via independence among the events for each $v$
%
\begin{align}\label{eq:randomdraw3}
    \Pro{\mathcal{E}_{\text{leaves}}} \geq \left(1- \frac{2t + 1}{n^\eps}\right)^{m} \in \Omega(1),
\end{align}
%
where we used that $t$ is constant and $m \in \bigO(n^{\eps/2})$.

Putting \Cref{eq:randomdraw1}, \Cref{eq:randomdraw2} and \Cref{eq:randomdraw3} together and using independence we then obtain
%
\begin{align}\label{eq:final-prob}
  \Pro{\mathcal{E}} = \Pro{\mathcal{E}_u \cap \mathcal{E}_{\text{path}} \cap \mathcal{E}_{\text{leaves}}} \in \Omega(n^{-(2t +1)\eps}).  
\end{align}
%
To finish the proof of our claim, let $X_u$ be the indicator random variable that event $\mathcal{E}$ occurs for vertex $u \in U_{(i)}$ and define the random variable $X := \sum_{u\in U_{(i)}} X_u \geq K(i+1)$. We then get by linearity of expectation $\Exp{X} \geq \Pro{\mathcal{E}}\cdot |U_{(i)}| \in \Omega\left(n^{1 - 1/2\alpha - \eps -(i+1)\eps(2t+1)} \right)$ by the assumption that $K_{(i)} \in  \Omega\left(n^{1 - 1/2\alpha - \eps -i\eps(2t+1)} \right)$ and \Cref{eq:final-prob}. Since the random variables $X_u$ are independent by the 'moreover' statement of \Cref{lem:luby-obstruction}, a Chernoff bound gives the desired result that $K_{(i+1)} \in \Omega\left(n^{1 - 1/2\alpha - \eps -(i+1)\eps(2t+1)} \right)$ \aas
\end{proof}}
%
Now, using that $K_{(0)} \in \Omega\left(n^{1 - 1/2\alpha - \frac{1 - \frac{1}{2\alpha}}{2t(t+2)}}\right)$ \aas by \Cref{lem:luby-obstruction}, and applying Claim~\ref{claim:appendix}iteratively for $t \in \bigO(1)$ rounds, a union bound over the $t$ induction steps yields
%
\begin{align*}
\Pro{K_{(t)} \in \Omega\left(n^{1 - 1/2\alpha - \eps -(t+1)\eps(2t+1)} \right)} &\geq 1 - \Pro{K_{(0)} \not \in \Omega\left(n^{1 - 1/2\alpha - \frac{1 - \frac{1}{2\alpha}}{2t(t+2)}}\right)}\\ 
&- \Pro{\bigcup_{i=0}^{t-1} K_{(i+1)} \not \in \Omega\left(n^{1 - 1/2\alpha - \eps -(i+1)\eps(2t+1)} \right) | K_{(i)}\in  \Omega\left(n^{1 - 1/2\alpha - \eps -i\eps(2t+1)} \right)}\\ &\in 1- o(1).
\end{align*}
%
This finishes the proof since by our choice $\eps := \frac{1 - \frac{1}{2\alpha}}{2t(t+2)} > 0$, $K_{(t)} \in n^{\Omega(1)}$ \aas
\end{proof}

\end{document}
